\section{Numerical Verification and Reproducibility}
\label{app:numerical-diagnostics}

This appendix records independent numerical checks of the quantities used
in the proof. They are designed to catch transcription, normalization, and
sign errors. They are not proof inputs: finite enumeration cannot cover all
dimensions, and a mesh cannot establish an inequality on a continuum. The
analytic arguments in the preceding appendices remain the certificates.

\begin{remark}[Reproducing the diagnostics]
From the source directory, run
\[
 \texttt{python3 numerical\_diagnostics.py}.
\]
The script uses NumPy for finite-cube enumeration and stable elementary
formulas for the scalar channel quantities. Information values are in
nats; scalar margins use the normalizations of the cited lemmas. Displayed
values are rounded only after the underlying expressions have been evaluated.
\end{remark}

\subsection{Exhaustive Checks on Small Cubes}

For a truth table \(f\) on \(n\) bits, the computation forms the full
noise matrix
\begin{align*}
 K_\rho(y,x)&=2^{-n}\prod_{i=1}^n(1+\rho x_i y_i),\\
 u(y)&=\sum_xK_\rho(y,x)f(x),
\end{align*}
and evaluates \(I_f(\rho)=h(\E f)-2^{-n}\sum_yh(u(y))\).
Thus there is no Monte Carlo error. We enumerated all \(2^{16}=65{,}536\)
Boolean truth tables for \(n=4\). We also enumerated all \(7{,}581\)
monotone truth tables for \(n=5\), generating them recursively from
lower and upper sections \(g\le h\).

At every tested channel the largest lead \(I_f-\phi\) was attained by
a coordinate rule and differed from zero by at most
\(3.4\times10^{-16}\). The table gives the positive deficit
\(\phi-I_f\) of the best rule after excluding dictators and
anti-dictators in dimension four, and after excluding coordinate rules
in the monotone dimension-five enumeration.
\begin{center}
{\small\setlength{\tabcolsep}{7pt}\renewcommand{\arraystretch}{1.1}
\begin{tabular}{@{}rrr@{}}
\toprule
\(\rho\) & all rules, \(n=4\) & monotone rules, \(n=5\)\\ \midrule
0.10 & \(8.70084\times10^{-4}\) & \(5.09022\times10^{-4}\)\\
0.25 & \(5.33512\times10^{-3}\) & \(3.15088\times10^{-3}\)\\
0.50 & \(1.94985\times10^{-2}\) & \(1.18769\times10^{-2}\)\\
0.75 & \(3.37645\times10^{-2}\) & \(2.12692\times10^{-2}\)\\
0.90 & \(3.11339\times10^{-2}\) & \(1.93117\times10^{-2}\)\\
0.99 & \(1.04605\times10^{-2}\) & \(6.06839\times10^{-3}\)\\
\bottomrule
\end{tabular}}
\end{center}
These deficits are not uniform as \(\rho\) approaches either endpoint;
that is expected and is one reason the proof uses an endpoint argument
rather than a global numerical gap.

\subsection{Low-Correlation Certificates}

The two tangent coefficients recompute as
\begin{align*}
 \alpha&=1.325831544044512\ldots,\\
 \beta&=0.849913438894783\ldots.
\end{align*}
Substitution into the expressions of Section~\ref{app:low-arithmetic}
gives the following values. The last column is the rational certificate
printed in the proof.
\begin{center}
{\small\setlength{\tabcolsep}{6pt}\renewcommand{\arraystretch}{1.08}
\begin{tabular}{@{}lrr@{}}
\toprule
Expression & Recomputed value & Certificate\\ \midrule
\(M_{\rm loc}(r_1,5/8)\) & \(8.014515\times10^{-4}\) & \(>1/1400\)\\
\(B_0'(5/8)\) & \(-6.848011\times10^{-2}\) & \(<-17/250\)\\
\(B_0''(5/8)\) & \(2.332248\times10^{-1}\) & \(<7/30\)\\
\(B_0'''(5/8)\) & \(-1.029815\) & \(<-1\)\\
\(B_0''(4/5)\) & \(-1.196782\times10^{-1}\) & \(<-1/9\)\\
\(B_0''(r_1)\) & \(-1.123197\times10^{-1}\) & \(<-1/10\)\\
\(B_0(r_1)\) & \(3.913927\times10^{-4}\) & \(>1/4000\)\\
Starting contraction slack & \(2.473064\times10^{-3}\) & \(>1/500\)\\
\bottomrule
\end{tabular}}
\end{center}
Every recomputed value lies on the strict side of its rational bound.
The narrowest displayed clearance is the curvature check at \(5/8\):
\(7/30-B_0''(5/8)=1.08491\times10^{-4}\).

\subsection{Intermediate-Range Reserves}

In this subsection, \(\alpha,\beta,\gamma,\lambda,\eta\) are the
middle-range parameters of Lemma~\ref{lem:middle-channel-checks};
\(\alpha\) and \(\beta\) no longer denote the tangent coefficients above.
For reference, define the half-height profile margin and the two
channel reserves by
\begin{align*}
 R_{\rm prof}(\ell)
 &=\log\frac{(\ell/2)F(\ell/2)}{\ell F(\ell)}
   -\beta\bigl[H(\ell)-H(\ell/2)\bigr],\\
 R_{\rm en}(\ell)&=\theta^2-rH(\ell)-\eta/2,\\
 R_{\rm mean}(\ell)
 &=\gamma-\frac{2\theta^2}{\sinh^2\ell}
                -\frac\lambda\alpha-r\ell.
\end{align*}
These are direct evaluations of the three comparisons in
Lemma~\ref{lem:middle-channel-checks}; the full profile comparison has
equality at \(v=\ell\), so \(R_{\rm prof}\) is shown at \(v=\ell/2\).
\begin{center}
{\small\setlength{\tabcolsep}{7pt}\renewcommand{\arraystretch}{1.08}
\begin{tabular}{@{}rrrr@{}}
\toprule
\(\ell\) & \(R_{\rm prof}\) & \(R_{\rm en}\) & \(R_{\rm mean}\)\\ \midrule
1.55 & 0.436092 & 0.752286 & 0.252406\\
2.00 & 0.607533 & 0.809600 & 0.593983\\
2.50 & 0.800374 & 0.691309 & 1.201045\\
3.00 & 1.001815 & 0.424392 & 2.248303\\
3.50 & 1.215187 & 0.057980 & 4.001328\\
\bottomrule
\end{tabular}}
\end{center}
A supplementary mesh used \(196\) channel heights at spacing \(0.01\)
on \([1.55,3.5]\) and the \(99\) ratios
\(v/\ell\in\{0.01,\ldots,0.99\}\). No negative profile margin was
found; the smallest sampled one was \(8.2344\times10^{-3}\), adjacent
to the equality endpoint. The smallest sampled energy and mean reserves
were respectively \(0.057980\) and \(0.252406\), matching the opposite
ends of the interval in the table.

\subsection{High-Correlation Margins}

The raw join and cutoff margins tend to zero as \(\ell\) grows. To
distinguish genuine loss of sign from the expected asymptotic scale, the
table reports \(e^{\ell/2}M_{\rm join}\) and
\(e^\ell M_{\rm cut}\), together with the unscaled small-height margin.
\begin{center}
{\small\setlength{\tabcolsep}{7pt}\renewcommand{\arraystretch}{1.08}
\begin{tabular}{@{}rrrr@{}}
\toprule
\(\ell\) & \(M_{\rm small}\) & \(e^{\ell/2}M_{\rm join}\) & \(e^\ell M_{\rm cut}\)\\ \midrule
3.5 & 0.158997 & 2.403324 & 4.548284\\
5 & 0.202903 & 3.896735 & 8.236764\\
10 & 0.253236 & 5.976483 & 18.337516\\
20 & 0.241820 & 6.279837 & 35.114182\\
\bottomrule
\end{tabular}}
\end{center}
On a \(0.01\)-spaced mesh of \(1{,}651\) heights from \(7/2\) to
\(20\), all three quantities stayed positive. Their smallest sampled
values were \(0.158997\), \(2.403324\), and \(4.548284\), all at the
shared endpoint \(\ell=7/2\). This agrees with the analytic proof's
choice to preserve the different asymptotic scales rather than seek a
uniform positive lower bound for the raw margins.

\paragraph{Outcome.}
The exhaustive enumerations found no finite-dimensional counterexample,
and the independently evaluated low, middle, and tail quantities showed
no sign or normalization discrepancy. This is useful corroboration, but
the theorem rests only on the analytic certificates printed earlier.
